% !TeX root = ../main.tex
\chapter{Aspetti di Gestione del Codice}

\section{Anatomia di un Progetto Java: Convenzioni e Struttura}

Nello sviluppo professionale, un progetto Java non è mai una semplice cartella piena di file sparsi. Esistono delle convenzioni universali (utilizzate da IDE come IntelliJ, Eclipse o sistemi come Maven e Gradle) che servono a separare ciò che il programmatore scrive da ciò che il computer esegue.

La \textbf{Root} (Radice) è la cartella principale che contiene l'intero progetto. Al suo interno, il lavoro viene diviso in due rami principali per garantire pulizia e sicurezza:

\begin{itemize}
    \item \textbf{Cartella \texttt{src} (Source)}: Contiene esclusivamente i file sorgente (\texttt{.java}). È l'unica cartella che viene salvata nei sistemi di controllo versione (come Git).
    \item \textbf{Cartella \texttt{bin} oppure \texttt{out} (Binary/Output)}:Qui il compilatore deposita i file \texttt{.class} (il bytecode). Questa cartella è generata automaticamente e può essere eliminata in qualsiasi momento senza perdere il lavoro, poiché viene rigenerata alla successiva compilazione.
\end{itemize}


\begin{observation}[Perché separare \texttt{src} da \texttt{bin}?]
    Separare i file è fondamentale per la distribuzione: se vuoi regalare il tuo programma a un amico, gli darai solo il contenuto di \texttt{bin}. Se gli dessi anche \texttt{src}, lui potrebbe vedere e modificare il tuo codice sorgente. Inoltre, permette a strumenti come Git di non tracciare file inutili che appesantirebbero solo il repository.
\end{observation}


\section{I Package e il Naming Convention}

Nello sviluppo software moderno, la probabilità che due programmatori diversi scelgano lo stesso nome per una classe (ad esempio \texttt{Employee}, \texttt{Account} o \texttt{User}) è estremamente alta. Per evitare conflitti, chiamati \textbf{collisioni di nomi} (\textit{namespace collisions}), Java introduce il concetto di package.

Un \textbf{package è una collezione di classi correlate che formano un'unità logica separata dal resto del codice}. Funge da "cognome" per la classe, garantendo che il suo nome sia unico all'interno dell'intero ecosistema Java.

Per garantire l'unicità dei package a livello globale, la comunità Java ha adottato una convenzione ingegnosa: utilizzare il proprio \textbf{nome di dominio Internet al contrario}.
Poiché i domini sono assegnati in modo univoco dalle autorità internazionali, nessuno potrà mai creare un package che inizi con il tuo dominio se tu ne sei il legittimo proprietario.

\begin{example}[Scomposizione di un nome]
    Se l'azienda \texttt{horstmann.com} sviluppa un progetto chiamato \texttt{corejava}, il package sarà \texttt{com.horstmann.corejava}. Una classe \texttt{Employee} al suo interno avrà come "nome legale completo" (\textbf{Fully Qualified Name}):
    \begin{center}
        \texttt{com.horstmann.corejava.Employee}
    \end{center}
\end{example}

\begin{warningbox}[Package vs Directory: Distinzione Concettuale]
Un errore tipico dei principianti è identificare il concetto di \textbf{package} con quello di \textbf{directory} (cartella). Sebbene Java, come vedremo a breve, imponga una corrispondenza biunivoca tra il nome del package e il percorso fisico sul disco, i due concetti rimangono distinti:
\begin{itemize}
    \item Il \textbf{Package} è un'entità logica (\textit{Namespace}): serve a raggruppare classi correlate e a prevenire conflitti di nomi.
    \item La \textbf{Directory} è un'entità fisica: è una struttura del \textit{file system} utilizzata dal sistema operativo.
\end{itemize}
In sintesi: un package \textbf{non è} semplicemente una collezione di file contenuti nella stessa cartella. Prova ne è il fatto che file situati nella stessa cartella, ma privi della direttiva \texttt{package}, appartengono, come vedremo, all'``unnamed package'' e non a un package che prende il nome dalla cartella stessa.
\end{warningbox}

Un altro errore comune per chi inizia è pensare che i package siano "annidati" come le cartelle di un sistema operativo e che quindi condividano permessi o caratteristiche. In Java, questa gerarchia è puramente visiva per l'essere umano, ma non esiste per il compilatore: per quest'ultimo, i package sono tutti "allo stesso livello". 

Non esiste dunque alcuna relazione tecnica tra package che sembrano annidati. Ad esempio, \texttt{java.util} e \texttt{java.util.jar} sono trattati come due entità \textbf{completamente distinte}. 
\begin{itemize}
    \item Una classe in \texttt{java.util} non ha alcun accesso privilegiato ai membri di una classe in \texttt{java.util.jar}.
    \item Importare \texttt{java.util.*} \textbf{non} importa le classi presenti in \texttt{java.util.jar}.
\end{itemize}

\begin{tip}[Best Practice per i Nomi]
    I nomi dei package dovrebbero essere sempre scritti in \textbf{minuscolo}. Questo aiuta a distinguerli immediatamente dalle classi, che per convenzione iniziano sempre con la lettera maiuscola.
\end{tip}

\begin{observation}[Titolarità dei Package]
    In ambito professionale, lo sviluppatore non utilizza un proprio dominio personale, bensì quello dell'organizzazione per cui lavora. 
    \begin{itemize}
        \item \textbf{Aziende}: Si utilizza il dominio aziendale (es. \texttt{com.google.*}).
        \item \textbf{Open Source}: Si utilizza spesso il prefisso del repository: \\[2pt] \quad\texttt{com.github.username.*}.
    \end{itemize}
    E importante notare che il compilatore Java \textbf{non verifica} l'esistenza reale del dominio; la convenzione serve esclusivamente a garantire l'unicità logica dei nomi su scala globale.
\end{observation}

\subsection{Collaborazione e Conflitti}

Quando più programmatori lavorano sullo stesso dominio aziendale, l'organizzazione segue regole di compartimentazione rigorose per evitare sovrapposizioni:

\begin{itemize}
    \item \textbf{Sotto-pacchetti Specifici}: Il progetto viene suddiviso in moduli funzionali (es. \texttt{ui}, \texttt{backend}, \texttt{security}), assegnando a ogni team o sviluppatore una porzione specifica dell'albero dei package.
    \item \textbf{Naming Convention del Ruolo}: Si prediligono nomi di classe che descrivano la funzione specifica (es. \texttt{UserValidator} invece di un generico \texttt{UserCheck}), riducendo il rischio di omonimie.
    \item \textbf{Sistemi di Controllo Versione}: Strumenti come \textbf{Git} gestiscono la presenza di più sviluppatori sugli stessi file, permettendo la fusione delle modifiche (\textit{merging}) e la risoluzione dei conflitti in modo controllato.
\end{itemize}

\begin{tip}[L'Importanza del Prefisso]
    Il prefisso del dominio (es. \texttt{it.ferrari}) serve a proteggere il progetto dal \textbf{mondo esterno}. L'organizzazione interna dei sub-package serve invece a proteggere i \textbf{colleghi tra loro}.
\end{tip}

\section{Creazione e Gestione Manuale dei Package}

Per comprendere davvero come Java gestisce il codice, è utile simulare il lavoro che l'IDE svolge "dietro le quinte". La creazione di un package senza l'aiuto di IntelliJ richiede una sincronizzazione precisa tra il codice sorgente e la struttura delle cartelle del sistema operativo.

\subsection*{Fase 1: Dichiarazione e Struttura Fisica}
Per creare un package, ad esempio \texttt{eu.dedalo.math}:
\begin{itemize}
  \item si crea un percorso fisico nella cartella \texttt{src} che rifletta il nome del package, dunque \texttt{src/eu/dedalo/math/},
  \item in tutti i file sorgenti .java che stanno alla fine del percorso si dichiara che appartengono a quel package: \texttt{package eu.dedalo.math;}. Questa direttiva si dice \textbf{dichiarazione di appartenenza al package (package statement)} e deve essere la prima istruzione del file, prima di qualsiasi import o dichiarazione di classe.
\end{itemize}

\vspace{\baselineskip}

In Java, la regola è che ogni file sorgente che dichiara di appartenere a un package deve essere posizionato in una cartella che rispecchia esattamente la struttura del package. Se questa corrispondenza non viene rispettata, il compilatore genererà errori. Se invece la dichiarazione di package è assente, il file viene considerato parte dell'``unnamed package'', che è una sorta di spazio dei nomi predefinito per i file senza package.


\subsection*{Fase 2: Compilazione Professionale}

Supponiamo ora di aver creato il file sorgente \texttt{Calcolatrice.java} nel package che abbiamo appena costruito. È tecnicamente possibile compilare il file puntando direttamente al suo percorso nel file system senza specificare una directory di output:

\begin{lstlisting}[language=bash]
# Compilazione diretta dal terminale
javac src/eu/dedalo/math/Calcolatrice.java
\end{lstlisting}

In questo caso, però, il compilatore non crea una struttura separata, ma deposita il file \texttt{.class} generato nella stessa cartella del file \texttt{.java}. 

Per questo motivo la best practice prevede l'uso del flag \texttt{-d} (\textit{destination}).

\begin{lstlisting}[language=bash, caption={Compilazione da terminale}]
# Posizionarsi nella cartella 'src'
javac -d ../bin eu/dedalo/math/Calcolatrice.java
\end{lstlisting}

Il comando \texttt{-d ../bin} istruisce il compilatore a:
\begin{itemize}
    \item Leggere la riga \texttt{package eu.dedalo.math;} nel sorgente.
    \item Creare automaticamente la gerarchia di cartelle \texttt{eu/dedalo/math/} dentro la cartella \texttt{bin}.
    \item Depositare il file \texttt{Calcolatrice.class} nella cartella finale.
\end{itemize}


\subsection*{Fase 3: Esecuzione e il "Nome Legale"}
L'errore più comune in questa fase è cercare di eseguire il programma entrando nella cartella dove risiede il file \texttt{.class}. In Java, questo non funziona.

Per lanciare il programma da terminale, la JVM deve essere posizionata sulla \textbf{base della gerarchia dei package} (comunemente chiamata \textit{Classpath Root}).

\begin{itemize}
    \item Se ti trovi fisicamente dentro la cartella \texttt{bin}, la base è la cartella corrente.
    \item Se ti trovi nella \textit{Project Root}, devi specificare dove sono i binari: \\ \texttt{java -cp bin eu.dedalo.math.Calcolatrice}
\end{itemize}

In ogni caso, \textbf{non} devi mai entrare dentro le sottocartelle del package (come \texttt{math/}) per lanciare il comando, altrimenti la JVM "perderebbe il filo" della struttura logica: la JVM cercherebbe una classe che dichiara di non avere package (unnamed package); trovando invece scritto \texttt{package eu.dedalo.math;}, la JVM rifiuterebbe di caricarla per incoerenza tra posizione e dichiarazione.


\begin{warningbox}[Compilazione ed Esecuzione: File vs Classi]
È fondamentale distinguere come il sistema si interfaccia con il codice a seconda della fase in cui ci troviamo. Assumendo di trovarsi nella \textbf{Project Root}:

\begin{itemize}
    \item \textbf{Il Compilatore (\texttt{javac})} lavora sui \textbf{File}:
    \begin{itemize}
        \item Ragiona per \textbf{percorsi fisici} (usa lo \textit{slash} \texttt{/}).
        \item Richiede l'estensione \textbf{\texttt{.java}}.
        \item \textit{Best Practice}: \texttt{javac -d bin src/eu/dedalo/math/Calcolatrice.java}
    \end{itemize}
    
    \item \textbf{La JVM (\texttt{java})} lavora sulle \textbf{Classi}:
    \begin{itemize}
        \item Ragiona per \textbf{nomi logici} (usa il \textit{dot} \texttt{.}).
        \item \textbf{Mai} usare l'estensione (niente \texttt{.class}).
        \item Richiede la \textbf{Classpath Root}: deve sapere da dove inizia il package (es. la cartella \texttt{bin}).
        \item \textit{Best Practice}: \texttt{java -cp bin eu.dedalo.math.Calcolatrice}
    \end{itemize}
\end{itemize}
\end{warningbox}

\begin{tip}[Regola Mnemonica]
    Usa lo \textbf{slash} (\texttt{/}) quando stai parlando con il Sistema Operativo (file sul disco). \\
    Usa il \textbf{punto} (\texttt{.}) quando stai parlando con la Java Virtual Machine (logica dei package).
\end{tip}

\section{L'Accesso a livello di Package}

Oltre ai modificatori \texttt{public} e \texttt{private}, Java prevede un terzo livello di visibilità che si attiva quando \textbf{non viene specificato alcun modificatore}. Questo livello è noto come \textit{package-private} o accesso di default.

Una classe, un metodo o una variabile definiti senza modificatori di accesso sono visibili e utilizzabili da \textbf{tutte le altre classi appartenenti allo stesso package}, ma sono totalmente inaccessibili all'esterno.

Sebbene l'accesso di package sia accettabile per le classi di utilità interna, Horstmann sottolinea come sia una scelta pericolosa per i campi di istanza (variabili).

\begin{warningbox}[Regola d'oro dell'Incapsulamento]
    Le variabili di istanza dovrebbero essere \textbf{sempre} dichiarate come \texttt{private}. Dimenticare il modificatore non le rende private, ma le espone a tutte le classi del package, violando il principio di protezione dei dati.
\end{warningbox}

\subsection{Sicurezza e Moduli}
In passato, era possibile "intrufolarsi" in un package altrui semplicemente dichiarando lo stesso nome di pacchetto nel proprio file sorgente. 
\begin{itemize}
    \item Il JDK oggi protegge i package di sistema (quelli che iniziano con \texttt{java.}).
    \item Per i progetti complessi, la soluzione moderna per isolare davvero i package è l'uso dei \textbf{Moduli} (introdotti in Java 9), che permettono di dichiarare esplicitamente quali package sono "esportati" e quali restano privati. Vedremo i moduli più avanti, ma è importante sapere fin da ora che rappresentano un'evoluzione significativa rispetto alla semplice organizzazione in package, offrendo un controllo molto più rigoroso sull'accesso e la visibilità del codice.
\end{itemize}

\section{Gli Import}

Una classe può utilizzare liberamente tutte le classi presenti nel proprio package e tutte le classi \textbf{public} provenienti da altri package. Per accedere a queste ultime, Java offre due strade: l'uso del nome completo o la scorciatoia dell'istruzione \texttt{import}.

Il modo più diretto (ma faticoso) per usare una classe di un altro package è scriverne il \textbf{Fully Qualified Name} ogni volta che la si invoca:
\begin{center}
    \texttt{java.time.LocalDate today = java.time.LocalDate.now();}
\end{center}
Poiché questo approccio è prolisso, si utilizza quasi sempre l'istruzione \texttt{import} in cima al file (sotto la dichiarazione del package). Questo agisce come un "alias" che permette di usare solo il nome semplice della classe.

\begin{observation}[Tipi di Import]
Esistono due modi per importare:
\begin{itemize}
    \item \textbf{Specifico}: \texttt{import java.time.LocalDate;} (Importa solo quella classe).
    \item \textbf{Wildcard}: \texttt{import java.time.*;} (Importa potenzialmente tutte le classi del package).
\end{itemize}
\textbf{Nota bene:} L'uso del carattere \texttt{*} non ha alcun impatto negativo sulla dimensione del file compilato (\textit{bytecode}), ma rende meno esplicito quali classi siano effettivamente utilizzate nel file.
\end{observation}

\begin{warningbox}[La regola del pacchetto singolo]
    La wildcard \texttt{*} può "aprire" solo un singolo package. Non è possibile usare \texttt{import java.*} per importare tutti i sotto-pacchetti che iniziano con \texttt{java}. Ogni livello va importato singolarmente.
\end{warningbox}

\begin{tip}[L'eccezione java.lang]
    Il package \texttt{java.lang} (che contiene classi base come \texttt{System}, \texttt{String} o \texttt{Math}) è \textbf{sempre} importato automaticamente in ogni file Java. Non c'è mai bisogno di scriverlo esplicitamente.
\end{tip}

\subsection{Risoluzione dei Conflitti di Omonimia}
Il problema principale dei package sorge quando due classi hanno lo stesso nome ma risiedono in pacchetti diversi. L'esempio classico è la classe \texttt{Date}, presente sia in \texttt{java.util} che in \texttt{java.sql}.

\begin{example}[Gestire i conflitti]
    Se importiamo entrambi i package con la wildcard:
    \begin{itemize}
        \item \texttt{import java.util.*;}
        \item \texttt{import java.sql.*;}
    \end{itemize}
    Il compilatore andrà in errore se scriviamo semplicemente \texttt{Date}, perché non saprà quale scegliere. Per risolvere:
    \begin{enumerate}
        \item Si aggiunge un import specifico per "rompere il pareggio": \\[2pt] \quad \texttt{import java.util.Date;}
        \item Se servono \textbf{entrambe} le classi nello stesso file, si è obbligati a usare il nome completo (\textit{Fully Qualified Name}) per almeno una delle due: \\
        \texttt{var d1 = new java.util.Date();} \\
        \texttt{var d2 = new java.sql.Date(...);}
    \end{enumerate}
\end{example}



\begin{observation}[Dietro le quinte del compilatore]
    Si ricordi che gli \texttt{import} sono solo una comodità per il programmatore. Una volta compilato il codice, nel \textbf{bytecode} finale non esistono più nomi brevi: il compilatore sostituisce ogni riferimento con il \textbf{Fully Qualified Name} completo.
\end{observation}

\subsection{Gli Import Statici}

Oltre alle classi, Java permette di importare direttamente \textbf{metodi e campi statici}. Una volta importati, questi membri possono essere utilizzati nel codice come se fossero definiti localmente, senza dover far precedere il nome della classe di appartenenza.

\begin{example}[Sintassi]
Per importare tutti i membri statici della classe \texttt{System}:
\begin{lstlisting}
import static java.lang.System.*;

out.println("Ciao!"); // Invece di System.out.println
exit(0);              // Invece di System.exit(0)
\end{lstlisting}
\end{example}

Horstmann solleva un punto critico: non sempre abbreviare è un bene.
\begin{itemize}
    \item \textbf{Sconsigliato}: Abbreviare \texttt{System.out} in \texttt{out} spesso rende il codice meno chiaro per chi legge, poiché \texttt{out} è un nome molto generico.
    \item \textbf{Consigliato}: L'uso con la classe \texttt{Math} è molto apprezzato. Scrivere:\\[2pt] \quad\texttt{sqrt(pow(x, 2))} \\[2pt] \quad è molto più vicino alla notazione matematica rispetto a: \\[2pt] \quad\texttt{Math.sqrt(Math.pow(x, 2))}.
\end{itemize}

\begin{tip}[Import di Enumerazioni]
Gli import statici sono estremamente comodi con le \textbf{Enum}. Se devi gestire molti giorni della settimana, importare \texttt{static java.time.DayOfWeek.*} ti permette di usare semplicemente \texttt{FRIDAY} invece di \texttt{DayOfWeek.FRIDAY}.
\end{tip}


\begin{observation}[Il bytecode non mente]
Proprio come per gli import di classe, gli import statici non hanno alcun impatto sulle prestazioni o sulla dimensione del file \texttt{.class}. Il compilatore si occupa semplicemente di "risolvere" i nomi abbreviati sostituendoli con i nomi completi durante la compilazione.
\end{observation}